\appendix

% ============================================================
\section{Pre-Registered Hypotheses}
\label{sec:appendix_preregistered}
% ============================================================

Table~\ref{tab:preregistered} records all 16 pre-registered predictions alongside
experimental outcomes. Predictions were added in five stages before running the
corresponding experiments.

\textbf{Hosting and criteria.}
The ledger is hosted in the project notebook
\texttt{RESEARCH\_JOURNEY\_DEFENSE.md} with stage-stamped entries; each stage
was committed before the next round of training or evaluation began. A
prediction is marked \emph{confirmed} ({\checkmark}) when its directional claim
\emph{and} any pre-stated quantitative threshold (e.g., ``$\geq 30\%$
reduction'') hold; \emph{refuted} ({\texttimes}) when either fails;
\emph{partial} when one holds and the other does not. Only secondary protocol
details---mask seed indices, reporting-horizon choice, and figure
formatting---were finalized post-hoc; the predictions themselves, including
quantitative thresholds and the model bank, were frozen at registration time.
The breakdown is ten confirmed, three refuted, and three partial; verbatim
re-statements are linked from each row of Table~\ref{tab:preregistered}.

\begin{table}[h]
\centering
\caption{Pre-registered hypotheses and outcomes. Confirmed ({\checkmark}) or refuted
(\texttimes) by experiment.}
\label{tab:preregistered}
\vskip 0.1in
\small
\begin{tabular}{p{6cm}p{5cm}c}
\toprule
\textbf{Prediction} & \textbf{Outcome} & \textbf{Result} \\
\midrule
GWNet fails hardest under city transfer (topology overfitting)
  & GWNet is worst under sensor dropout, not city transfer; GRU is worst at city transfer
  & \texttimes \\
Graph-based models are more robust than GRU to sensor dropout
  & GRU is nearly as robust as DCRNN at 30\%; adaptive graph hurts, fixed graph helps
  & Partial \texttimes \\
All models degrade on PEMS-BAY vs.\ METR-LA
  & Source-stat MAE improves on PEMS-BAY (simpler traffic); target-stat reveals real gaps
  & \texttimes \\
Adaptive adjacency drives sensor-dropout fragility ($\geq$30\% reduction when ablated)
  & 43.1\% reduction confirmed
  & {\checkmark} \\
Temporal shift magnitude differs across architectures
  & All architectures within one std of each other
  & \texttimes \\
STGCN (fixed spectral adj) $\DeltaSensor \leq$ DCRNN (fixed diffusion adj)
  & STGCN $+0.078$ $>$ DCRNN $+0.020$; but both $<$ GWNet $+0.088$. Adaptive adj is the dominant driver.
  & Partial \texttimes \\
\midrule
\multicolumn{3}{l}{\emph{P11--P13 added prior to structured / coverage experiments:}} \\
{[P11]} GWNet transfer gap persists (GWNet worst) across 150- and 207-node PEMS-BAY coverage
  & Confirmed: ranking preserved; GWNet +41\% MAE at 150 nodes vs.\ $-2\%$ for DCRNN
  & {\checkmark} \\
{[P12]} Hub dropout raises GWNet $\DeltaSensor$ relative to uniform random (dense routing)
  & Confirmed: GWNet $+38\%$ ($+0.127$); DCRNN lower ($+0.008$ vs.\ random $+0.021$)
  & {\checkmark} \\
{[P13]} Mask-aware recovery meaningful with FP $\leq5\%$, FN $\leq10\%$ detection
  & Partial: 27\% recovery at FP=5\%/FN=10\%; high variance at FP=10\%/FN=20\%
  & Partial {\checkmark} \\
\midrule
\multicolumn{3}{l}{\emph{P14--P16 added prior to inductive-baseline and capacity-matched experiments:}} \\
{[P14]} MTGNNLight (sparse cosine-sim) shows lower $\DeltaSensor$ than GWNet (dense product)
  & Confirmed: MTGNNLight $+0.051$, GWNet $+0.088$ ($42\%$ lower)
  & {\checkmark} \\
{[P15]} GWNet[CM] (doubled patience) shows same $\DeltaSensor$ as standard GWNet
  & Confirmed: $+0.088$ vs $+0.088$ (identical); fragility is architectural not training artifact
  & {\checkmark} \\
{[P16]} Source-only adj control narrows but does not eliminate DCRNN's transfer advantage
  & Confirmed: DCRNN $0.180$ vs GWNet $0.268$ ($49\%$ gap remains under fair comparison)
  & {\checkmark} \\
\midrule
\multicolumn{3}{l}{\emph{P17--P19 added prior to stronger-GWNet, second-city-pair, and per-horizon experiments:}} \\
{[P17]} GWNet[Strong] (hidden 128, embed 20) still shows $\DeltaSensor{>}2{\times}$ DCRNN and unchanged city MAE
  & Confirmed: $+0.050$ vs $+0.020$ ($2.5{\times}$); city $4.46$ vs $4.42$ mph
  & {\checkmark} \\
{[P18]} DCRNN is best-generalizing on $\geq1$ PEMS04-source direction (PEMS04$\to$METR-LA, PEMS04$\to$PEMS-BAY)
  & Confirmed on \emph{both}: DCRNN $7.01$/$2.65$ vs next-best GRU $11.03$/GWNet $4.58$
  & {\checkmark} \\
{[P19]} GWNet's per-horizon city-transfer MAE is flatter than DCRNN's (static mis-routing, not accumulation)
  & Confirmed: GWNet $+4\%$ h1$\to$h12 growth vs DCRNN $+46\%$
  & {\checkmark} \\
\midrule
\multicolumn{3}{l}{\emph{P21 added prior to full-sweep best-config experiment:}} \\
{[P21]} GWNet[Best] (12-config sweep: hidden 256, embed 32, lr 5e-4) still shows $\DeltaSensor{>}2{\times}$ DCRNN
  & Confirmed: $+0.042$ vs $+0.020$ ($2.1{\times}$); in-dist $7.22$ mph; optimal training cannot eliminate routing fragility
  & {\checkmark} \\
\bottomrule
\end{tabular}
\end{table}

% ============================================================
\section{Protocol and Transparency}
\label{sec:appendix_protocol_group}
% ============================================================

\subsection{Estimand-level crosswalk}
\label{sec:appendix_estimand_crosswalk}

Table~\ref{tab:estimand_crosswalk} maps each reported estimand to its source
table, units, normalization convention, and whether it is comparable across
model families. Clean MAE is reported in canonical (mph) form for the primary
table and in each family's standard masked convention for the extended bank;
clean values are not directly cross-family-rankable because per-family
training conventions differ. Shift estimands ($\DeltaSensor$, cross-city,
$\DeltaTemp$) are normalized so that within-family and (where matched
structural access holds) across-family comparisons are sound.

\begin{table}[h]
\centering
\caption{Estimand-level crosswalk: which numbers are cross-family-comparable.}
\label{tab:estimand_crosswalk}
\vskip 0.05in
\small
\resizebox{\linewidth}{!}{%
\begin{tabular}{lllll}
\toprule
\textbf{Estimand} & \textbf{Source} & \textbf{Unit} & \textbf{Normalization} & \textbf{Cross-family?} \\
\midrule
Clean MAE              & T.\ref{tab:main}, T.\ref{tab:extended}                              & mph (canon.) / masked (ext.) & per-family             & No \\
$\DeltaSensor$         & T.\ref{tab:main}, T.\ref{tab:extended}                              & normalized                   & source-stat            & Yes (within-family) \\
Cross-city MAE         & T.\ref{tab:main}, T.\ref{tab:extended}                              & normalized                   & source / target (both) & Yes (matched-access) \\
$\DeltaTemp$           & T.\ref{tab:main}, T.\ref{tab:extended}                              & normalized                   & source-stat            & Yes \\
Routing concentration  & T.\ref{tab:routing_spectrum}, App.~\ref{sec:appendix_pems04_routing} & unitless $[0,1]$            & per-row                & Yes (proxy-matched, $n{=}4$) \\
\bottomrule
\end{tabular}}
\end{table}

The ``Cross-family?'' column reflects whether the estimand can be ranked
across model families without re-introducing a convention artifact: clean MAE
cannot, the three shift estimands can (after the same source-stat
normalization), and routing concentration can only across the four
row-stochastic backbones for which the proxy is identically defined.

\subsection{Reproducibility, assets, and compute}
\label{sec:appendix_reproducibility}

All experiments use public traffic datasets and frozen checkpoints. The main
numeric source for Tables~\ref{tab:main}--\ref{tab:extended} is
\texttt{results/official\_adaptive\_relation/}\allowbreak\texttt{UPDATED\_MERGED\_COMPARISON\_TABLE\_2026-04-16.csv};
PEMS04 routing replication and mask-aware sweeps are read from
\texttt{results/polish\_v15/}. The checkpoint registry is
\texttt{baselines/stgnn\_protocol\_registry.json}, and the unified training
configuration is \texttt{configs/stgnn\_unified.yaml}. Evaluation commands are
listed in the repository README and mirror the commands used to generate the
tables and figures.

The released artifact contains the recommendation table and the scripts needed
to regenerate the shift profiles from the canonical result files. It does not
introduce a new dataset. METR-LA, PEMS-BAY, and PEMS04 remain under their
original distribution terms; we cite the source papers and do not repackage the
raw datasets in the submission. The runs were executed with one CUDA GPU per
training/evaluation job. The April 16 extended-model evaluation log records
roughly 3--9 minutes per seed for shift evaluation after training, depending on
backbone; full retraining is not required to verify the frozen-checkpoint
results.

% ============================================================
\section{Sensor-Dropout Analysis}
\label{sec:appendix_sensor_group}
% ============================================================

\subsection{Severity curve: dropout ratios 10--50\%}
\label{sec:appendix_severity}

Table~\ref{tab:sweep} reports $\DeltaSensor$ for three backbones across dropout
ratios from 10\% to 50\%.

\begin{table}[h]
\centering
\caption{\textbf{GWNet is the most fragile backbone at every dropout level; the gap grows
nonlinearly.} $\DeltaSensor$ mean over 3 seeds across dropout ratios 10--50\%.
GWNet/DCRNN ratio (bottom) quantifies relative fragility.}
\label{tab:sweep}
\vskip 0.05in
\begin{tabular}{lccccc}
\toprule
\textbf{Backbone} & \textbf{10\%} & \textbf{20\%} & \textbf{30\%} & \textbf{40\%} & \textbf{50\%} \\
\midrule
DCRNN & $+0.005$ & $+0.012$ & $+0.020$ & $+0.033$ & $+0.044$ \\
GWNet & $+0.013$ & $+0.044$ & $+0.088$ & $+0.149$ & $+0.189$ \\
GRU   & $+0.001$ & $+0.005$ & $+0.025$ & $+0.073$ & $+0.088$ \\
\midrule
GWNet/DCRNN & $2.5\times$ & $3.8\times$ & $4.3\times$ & $4.6\times$ & $4.3\times$ \\
\bottomrule
\end{tabular}
\end{table}

\subsection{Multi-mask aggregation}
\label{sec:appendix_multimask}

Table~\ref{tab:multimask} reports $\DeltaSensor$ averaged over 5 random 30\%
zero-masks (seeds 42--46) per model checkpoint.
The main-text results use the pre-registered primary mask per model seed.
For DCRNN, MTGNNLight, STGCN, and GWNet the 5-mask aggregate remains close to
Table~\ref{tab:main}; GRU is a mask-sampling caveat.

\begin{table}[h]
\centering
\caption{Multi-mask sensor dropout: $\DeltaSensor$ mean and std over 5 random masks
(seeds 42--46), then averaged over 3 model seeds.
Most values are consistent with primary-mask results in Table~\ref{tab:main};
GRU is the exception noted in the protocol.}
\label{tab:multimask}
\vskip 0.1in
\resizebox{\linewidth}{!}{%
\begin{tabular}{lccc}
\toprule
\textbf{Backbone} & \textbf{$\DeltaSensor$ (multi-mask mean)} & \textbf{Seed std} & \textbf{vs.\ Table~\ref{tab:main}} \\
\midrule
DCRNN (fixed diffusion) & $+0.021$ & $\pm 0.002$ & $+0.020$ (consistent) \\
GRU   (graph-free)      & $+0.044$ & $\pm 0.002$ & $+0.025$ (seed-42 was low-variance mask) \\
MTGNNLight (sparse cosine-sim) & $+0.051$ & $\pm 0.012$ & $+0.051$ (consistent) \\
STGCN (fixed spectral)  & $+0.079$ & $\pm 0.006$ & $+0.078$ (consistent) \\
GWNet (adaptive adj)    & $+0.092$ & $\pm 0.030$ & $+0.088$ (consistent; high mask variance) \\
GWNet[CM] (patience=30) & $+0.088$ & $\pm 0.017$ & $+0.088$ (P15 confirmed identical) \\
\bottomrule
\end{tabular}}
\end{table}

\subsection{Structured dropout patterns}
\label{sec:appendix_structured_dropout}

Table~\ref{tab:structured_dropout} reports $\DeltaSensor$ under two topology-aware
dropout patterns alongside the uniform random baseline.
\emph{Hub dropout} removes the $k$ sensors with the highest adjacency-degree (most
message-passing traffic); \emph{contiguous dropout} removes a BFS-connected cluster
seeded from the highest-degree node, mimicking a geographic outage.

\begin{table}[h]
\centering
\caption{Structured vs.\ uniform sensor dropout ($\DeltaSensor$, mean over 3 seeds, 30\% drop rate).
Hub: top-$k$ degree sensors dropped.
Contiguous: BFS cluster from highest-degree seed.
GWNet ordering is preserved ($>$STGCN$>$DCRNN); structured patterns amplify GWNet's fragility
but reduce DCRNN/STGCN's apparent sensitivity, because hub dropout keeps peripheral
(easier-to-predict) sensors.}
\label{tab:structured_dropout}
\vskip 0.1in
\small
\begin{tabular}{lccc}
\toprule
\textbf{Backbone} & \textbf{Random (uniform)} & \textbf{Hub dropout} & \textbf{Contiguous dropout} \\
\midrule
DCRNN (fixed diffusion) & $+0.021$ & $+0.008$ & $+0.005$ \\
STGCN (fixed spectral)  & $+0.079$ & $+0.026$ & $+0.022$ \\
GWNet (adaptive adj)    & $+0.088$ & $+0.127$ & $+0.096$ \\
\bottomrule
\end{tabular}
\end{table}

The relative ordering GWNet $>$ STGCN $>$ DCRNN is preserved under all patterns.
GWNet's hub-dropout sensitivity is \textbf{45\% higher} than uniform random ($+0.127$ vs.\ $+0.088$,
matched single-mask protocol shared with Table~\ref{tab:main}):
hub nodes feed GWNet's densest adaptive-routing pathways, so removing them is especially damaging.
DCRNN and STGCN show \emph{lower} $\DeltaSensor$ under hub/contiguous dropout ($+0.008$, $+0.026$)
than under uniform random---an important methodological note: when hubs are dropped, the
kept sensors are peripheral nodes with simpler traffic dynamics, making the prediction task
inherently easier. The clean-kept and masked-kept MAEs both drop for DCRNN/STGCN under hub
dropout, but the \emph{gap} also shrinks, since DCRNN's bidirectional diffusion redistributes
information effectively even when hub inputs are absent.

\subsection{Node-degree stratification}
\label{sec:appendix_degree_stratification}

Table~\ref{tab:degree_stratification} reports $\DeltaSensor$ stratified by node degree
quartile (Q1 = lowest degree, Q4 = highest).
Nodes in each quartile are defined by their adjacency row-sum degree in the METR-LA graph.
Only kept sensors (not dropped) are included in the MAE computation.
Reported values are mean over 3 model seeds.

\begin{table}[h]
\centering
\caption{Sensor-dropout $\DeltaSensor$ by node-degree quartile (30\% zero-mask, kept sensors only, 3-seed mean). Mean node degrees: Q1\,$\bar{d}=1.83$, Q2\,$\bar{d}=3.32$, Q3\,$\bar{d}=4.48$, Q4\,$\bar{d}=5.98$.}
\label{tab:degree_stratification}
\vskip 0.1in
\small
\begin{tabular}{lccccl}
\toprule
\textbf{Backbone} & \textbf{Q1 (low)} & \textbf{Q2} & \textbf{Q3} & \textbf{Q4 (high)} & \textbf{Q4/Q1} \\
\midrule
DCRNN           & $+0.019$ & $+0.026$ & $+0.018$ & $+0.018$ & $0.95\times$ \\
MTGNNLight      & $+0.044$ & $+0.052$ & $+0.050$ & $+0.055$ & $1.25\times$ \\
STGCN           & $+0.059$ & $+0.089$ & $+0.074$ & $+0.089$ & $1.50\times$ \\
GWNet           & $+0.066$ & $+0.076$ & $+0.085$ & $+0.121$ & $1.83\times$ \\
\bottomrule
\end{tabular}
\end{table}

GWNet is the only backbone with a strongly monotonic degree$\to$fragility pattern.
MTGNNLight's weak positive relationship ($1.25\times$) and STGCN's intermediate ratio ($1.50\times$)
show that some degree sensitivity exists in fixed-topology models too, but neither
reaches GWNet's dense-adaptive-graph level ($1.83\times$), and only GWNet's increase is monotonic across
Q1$\to$Q4.

\subsection{Mask-aware recovery: FP/FN sensitivity}
\label{sec:appendix_mask_fpfn}

The mask-aware adjacency variant zeroes out adjacency rows/columns for detected failed sensors.
In silent-failure scenarios, the detector is imperfect: false positives (healthy sensors
flagged as failed) and false negatives (failed sensors missed) can both occur.
Table~\ref{tab:mask_fpfn} reports MAE under the perfect mask vs.\ two realistic detector
operating points.

\begin{table}[h]
\centering
\caption{Mask-aware GWNet sensitivity to imperfect mask detection (3 model seeds).
FP = fraction of healthy sensors erroneously zeroed in the adjacency;
FN = fraction of failed sensors left un-zeroed.
$\DeltaSensor$: estimated as $\DeltaSensor^{\text{perfect}} + \delta_{\text{from perfect}}$
where $\DeltaSensor^{\text{perfect}} = +0.023$ (from Table~\ref{tab:mask_aware}) and
$\delta_{\text{from perfect}}$ is the mean degradation relative to perfect detection.}
\label{tab:mask_fpfn}
\vskip 0.1in
\small
\begin{tabular}{lccc}
\toprule
\textbf{Detector scenario} & \textbf{FP / FN} & \textbf{$\delta_{\text{from perfect}}$ (mean)} & \textbf{Recovery vs.\ no mask-aware} \\
\midrule
Standard (no mask-aware) & ---       & ---    & 0\%   \\
Perfect mask             & 0\% / 0\% & $0.000$ & 75\%  \\
Conservative detector    & 5\% / 10\% & $+0.044$ & $\sim$27\% \\
Noisy detector           & 10\% / 20\% & $+0.077^*$ & $\sim$17\% \\
\bottomrule
\end{tabular}
\end{table}
\noindent\small{$^*$High seed variance at 10\%/20\% (seeds range: $+0.049$--$+0.134$); median is $+0.049$.}

The mask-aware approach retains meaningful benefit under realistic detector quality
(FP $\leq5\%$, FN $\leq10\%$), though the recovery drops from 75\% (perfect) to $\sim$27\%
at the conservative operating point---a reported result that motivates improving detector
precision before deploying the intervention.
At 10\%/20\% error rates, high variance (seed-dependent behaviour) makes the benefit
uncertain; we report the median recovery ($\sim$22\%) as more representative.
Sustained zero-reading anomaly detectors typically achieve FP $<5\%$ and FN $<10\%$
on traffic sensors, making the conservative regime ($\sim$27\% recovery) the practical target.

\subsection{Mask-aware dropout-ratio sweep}
\label{sec:appendix_mask_aware_sweep}

The mask-aware GWNet variant (Section~\ref{sec:mechanism}) is reported in
the main text at a single 30\% random-mask operating point. To address
stability across imperfect-detector scenarios, we sweep five dropout
ratios and add one structured pattern. All cells use 3 GWNet seeds
$\times$ 5 mask seeds; structured pattern uses the top-30\% nodes by
adjacency degree (one mask per seed).

\begin{table}[h]
\centering
\small
\caption{Mask-aware GWNet recovery is stable across dropout ratios and
structured patterns. ``Standard'' = silent zero-mask; ``mask-aware'' =
adjacency rows/cols zeroed and renormalized for dropped sensors;
recovery $= 1 - \DeltaSensor^{\text{mask-aware}}/\DeltaSensor^{\text{standard}}$.
Script: \texttt{experiments/mask\_aware\_dropout\_sweep.py};
data: \texttt{results/polish\_v15/mask\_aware\_ratio\_sweep.json}.}
\label{tab:mask_aware_sweep}
\vskip 0.05in
\begin{tabular}{lccc}
\toprule
\textbf{Cell} & \textbf{$\DeltaSensor^{\text{std.}}$} & \textbf{$\DeltaSensor^{\text{mask-aware}}$} & \textbf{Recovery} \\
\midrule
Random $5\%$        & $+0.004 \pm 0.002$ & $+0.000 \pm 0.001$ & $109\% \pm 29$ \\
Random $10\%$       & $+0.019 \pm 0.006$ & $+0.005 \pm 0.002$ & $\phantom{0}76\% \pm \phantom{0}6$ \\
Random $20\%$       & $+0.051 \pm 0.013$ & $+0.015 \pm 0.002$ & $\phantom{0}70\% \pm \phantom{0}5$ \\
Random $30\%$ (main) & $+0.092 \pm 0.030$ & $+0.023 \pm 0.009$ & $\phantom{0}75\% \pm \phantom{0}3$ \\
Random $40\%$       & $+0.145 \pm 0.026$ & $+0.035 \pm 0.006$ & $\phantom{0}75\% \pm \phantom{0}3$ \\
\midrule
Hub-30\% (structured) & $+0.116 \pm 0.037$ & $+0.034 \pm 0.001$ & $\phantom{0}71\%$ \\
\bottomrule
\end{tabular}
\end{table}

The mask-aware recovery is in the $70$--$76\%$ band across the full sweep
(the $5\%$ row's nominal $109\%$ reflects an absolute $\DeltaSensor^{\text{mask-aware}}$
within seed noise of zero). Hub-dropout, where 30\% of \emph{high-degree}
nodes fail simultaneously, is recovered at $71\%$---comparable to the
random-30\% main-text result---supporting the argument that the mechanism
generalizes beyond the random-uniform mask used in the headline claim.

\subsection{LOCF and mean imputation baselines}
\label{sec:appendix_imputation}

Table~\ref{tab:imputation} compares three imputation strategies for dropped sensors:
\emph{zero-mask} (set dropped sensor channels to 0), \emph{LOCF} (last observed value,
i.e., the last input timestep value is carried forward for the entire window), and
\emph{mean} (per-node mean computed over the test split).

\begin{table}[h]
\centering
\caption{Sensor-dropout $\Delta$ under three imputation strategies (30\% drop rate, seed=42, 3 model seeds). LOCF improvement: fraction of zero-mask gap closed by LOCF. Negative improvement indicates worse than zero-mask.}
\label{tab:imputation}
\vskip 0.1in
\small
\begin{tabular}{lcccc}
\toprule
\textbf{Backbone} & \textbf{$\Delta_{\text{zero}}$} & \textbf{$\Delta_{\text{LOCF}}$} & \textbf{$\Delta_{\text{mean}}$} & \textbf{LOCF improvement} \\
\midrule
DCRNN      & $+0.020$ & $\approx 0$ & $+0.025$ & $\sim 102\%$ \\
MTGNNLight & $+0.051$ & $\approx 0$ & $+0.077$ & $\sim 100\%$ \\
STGCN      & $+0.078$ & $+0.003$   & $+0.079$ & $\sim 97\%$ \\
GWNet      & $+0.088$ & $\approx 0$ & $+0.109$ & $\sim 101\%$ \\
\bottomrule
\end{tabular}
\end{table}

LOCF essentially eliminates the zero-mask degradation for all four backbones.
Mean imputation performs similarly to or worse than zero-masking, because the test-split
mean is influenced by the training distribution and creates a biased proxy in shifted
traffic conditions.
The LOCF result shows that the architectural fragility ranking we measure (GWNet
$>$ STGCN $>$ DCRNN) is specific to the zero-masking failure mode; any system with
temporal carry-forward (including simple telemetry holdover) would largely eliminate
the architectural gap.

% ============================================================
\section{Cross-City Transfer Analysis}
\label{sec:appendix_city_group}
% ============================================================

\subsection{Normalization conventions}
\label{sec:appendix_protocol}

Table~\ref{tab:norm_variants} reports cross-city evaluation under both normalization
conventions.
\emph{Source-stat} normalizes and denormalizes with METR-LA training statistics;
\emph{target-stat} uses PEMS-BAY training statistics for denormalization only
(PEMS-BAY statistics are not available at deployment time, so source-stat is the
operationally correct zero-shot metric; target-stat is reported for comparability
with the broader traffic forecasting literature).

\begin{table}[h]
\centering
\caption{Cross-city MAE under two normalization conventions.
Source-stat: METR-LA train mean/std for both normalization and denormalization
(lower = PEMS-BAY traffic is statistically simpler in normalized space).
Target-stat: PEMS-BAY train statistics for denormalization (interpretable mph units).
Mean $\pm$ population std over 5 random 207-node subsets and 3 model seeds.}
\label{tab:norm_variants}
\vskip 0.1in
\resizebox{\linewidth}{!}{%
\begin{tabular}{lccc}
\toprule
\textbf{Backbone} & \textbf{Clean MAE (METR-LA, norm)} & \textbf{Source-stat MAE (norm)} & \textbf{Target-stat MAE (mph)} \\
\midrule
DCRNN (fixed diffusion) & $0.359$ & $0.168$ & $2.74 \pm 0.01$ \\
STGCN (fixed spectral)  & $0.388$ & $0.209$ & $3.27 \pm 0.02$ \\
GWNet (adaptive adj)    & $0.427$ & $0.268$ & $4.42 \pm 0.15$ \\
GRU   (graph-free)      & $0.408$ & $0.278$ & $4.73 \pm 0.19$ \\
\bottomrule
\end{tabular}}
\end{table}

The source-stat improvement (PEMS-BAY $<$ METR-LA in normalized space) for all
backbones reflects lower traffic variance on PEMS-BAY highway sensors, not better
generalization. Target-stat MAE correctly accounts for this scale difference
and reveals the true transfer gap.

\subsection{Subset-size sensitivity}
\label{sec:appendix_subset_sensitivity}

To test whether the GWNet transfer gap depends on how many PEMS-BAY sensors are used,
we evaluate with 150-node coverage: 150 PEMS-BAY sensors are selected and the remaining
57 input slots are zero-padded (same protocol as sensor dropout).
Table~\ref{tab:subset_sensitivity} reports city transfer MAE under both coverages.

\begin{table}[h]
\centering
\caption{Cross-city transfer MAE (target-stat, mph) under variable sensor coverage.
150-node: 150 PEMS-BAY sensors active, 57 input slots zero-padded to 207.
207-node: full model capacity (different subset draws from main table; see text).
Mean $\pm$ std over 3 random PEMS-BAY subsets, 3 model seeds.}
\label{tab:subset_sensitivity}
\vskip 0.1in
\small
\begin{tabular}{lcc}
\toprule
\textbf{Backbone} & \textbf{150-node coverage} & \textbf{207-node (full)} \\
\midrule
DCRNN (fixed diffusion) & $4.25 \pm 0.03$ & $4.33 \pm 0.05$ \\
STGCN (fixed spectral)  & $4.94 \pm 0.25$ & $5.07 \pm 0.33$ \\
GWNet (adaptive adj)    & $7.49 \pm 1.08$ & $5.30 \pm 0.25$ \\
\bottomrule
\end{tabular}
\end{table}

The backbone ranking DCRNN $<$ STGCN $<$ GWNet is preserved at both coverage levels.
GWNet's MAE increases by \textbf{41\%} when moving from full 207-node to 150-node coverage
($5.30 \to 7.49$ mph); DCRNN decreases marginally ($-2\%$) and STGCN decreases marginally
($-3\%$).
The large GWNet penalty arises because its source-city node embeddings route information
through 57 zero-padded positions: inactive input slots inject systematic zero signals into
the adaptive adjacency pathways, compounding the city-transfer gap.
DCRNN and STGCN with target-city subgraph adjacency are unaffected by coverage reduction
because zero-padded nodes are disconnected in the graph.
Note: absolute MAE values differ from Table~\ref{tab:main} due to different random subset
draws; rankings are consistent.

\subsection{Source-only vs.\ target adjacency}
\label{sec:appendix_source_adj}

The main Table~\ref{tab:main} uses a \emph{matched structural constraint} (source-only
adjacency) as the canonical cross-city comparison, so that no backbone benefits from
access to target-city topology that others lack.
This appendix reports the full comparison: source-only normalized MAE (main table),
normal target-adj normalized MAE, and the original target-adj \emph{mph} values for
readers who prefer absolute units.
GWNet and MTGNNLight are source-tied by construction (they learn or compute their own
adjacency from source node embeddings); GRU uses no adjacency.

\begin{table}[h]
\centering
\caption{Full cross-city transfer comparison.
\emph{Source-only (norm)}: all models constrained to METR-LA adjacency; canonical main-table values.
\emph{Target-adj (norm)}: DCRNN/STGCN use PEMS-BAY adjacency (GWNet/MTGNNLight/GRU unchanged).
\emph{Target-adj (mph)}: same target-adj evaluation, denormalized using PEMS-BAY statistics.
Mean $\pm$ population std over 5 PEMS-BAY subsets, 3 model seeds.}
\label{tab:source_adj_control}
\vskip 0.1in
\small
\begin{tabular}{lccc}
\toprule
\textbf{Backbone} & \textbf{Source-only (norm)} & \textbf{Target-adj (norm)} & \textbf{Target-adj (mph)} \\
 & \textbf{[main table]} & & \\
\midrule
DCRNN (fixed diffusion) & $0.180 \pm 0.001$ & $0.167 \pm 0.002$ & $2.74 \pm 0.01$ \\
MTGNNLight (cosine-sim) & $0.244 \pm 0.004$ & $0.244 \pm 0.004$ & $4.15 \pm 0.11$ \\
STGCN (fixed spectral)  & $0.249 \pm 0.008$ & $0.202 \pm 0.005$ & $3.27 \pm 0.02$ \\
GRU   (graph-free)      & $0.278 \pm 0.009$ & $0.278 \pm 0.009$ & $4.73 \pm 0.19$ \\
GWNet (dense adaptive)  & $0.268 \pm 0.006$ & $0.268 \pm 0.006$ & $4.42 \pm 0.15$ \\
\bottomrule
\end{tabular}
\end{table}

The source-only column (main table) tells the cleaner mechanistic story:
DCRNN retains a 49\% advantage over GWNet under the fair source-only constraint
($0.180$ vs $0.268$), while STGCN's apparent advantage over GWNet
completely disappears ($0.249 \approx 0.268$).
STGCN's city-transfer advantage in the target-adj column ($0.202$) arises
from access to the target-city adjacency matrix, not from architectural generalization.
DCRNN's bidirectional diffusion generalizes better even when confined to source topology,
supporting its designation as the most transfer-robust backbone in this study.

\subsection{PEMS04-source second city pair (P18)}
\label{sec:appendix_pems04}

To test whether DCRNN's transfer advantage is METR-LA-specific, we train all four
backbones (DCRNN, GWNet, GWNet[Strong], GRU) on PEMS04~\citep{guo2019astgcn}
(307 sensors, California District 4 freeway, 16{,}992 timesteps) with the same frozen
training protocol and evaluate zero-shot on METR-LA and PEMS-BAY, using 5 random
307-node target subsets (seeds 42--46). The source PEMS04 adjacency is used for all
backbones at inference, matching the P16 source-only fairness control.

\begin{table}[h]
\centering
\caption{PEMS04-source cross-city transfer (3 seeds $\times$ 5 target subsets, mean $\pm$ std).
PEMS04 in-dist MAE in native flow units; target MAE in target-stat mph.
DCRNN is best in \emph{both} new directions by $>37\%$ margin.}
\label{tab:pems04_cross_city}
\vskip 0.1in
\small
\begin{tabular}{lccc}
\toprule
\textbf{Backbone} & \textbf{PEMS04 in-dist} & \textbf{$\to$ METR-LA (mph)} & \textbf{$\to$ PEMS-BAY (mph)} \\
\midrule
DCRNN         & $26.19 \pm 0.02$ & $\mathbf{7.01 \pm 0.09}$ & $\mathbf{2.65 \pm 0.02}$ \\
GWNet         & $33.47 \pm 0.45$ & $11.39 \pm 0.62$ & $4.58 \pm 0.16$ \\
GWNet[Strong] & $32.62 \pm 0.17$ & $12.58 \pm 0.32$ & $4.61 \pm 0.02$ \\
GRU           & $25.81 \pm 0.06$ & $11.03 \pm 0.01$ & $4.68 \pm 0.03$ \\
\bottomrule
\end{tabular}
\end{table}

DCRNN remains the best-generalizing backbone in both PEMS04-source directions, mirroring
the METR-LA-source and the PEMS-BAY $\to$ METR-LA reverse-transfer finding in
Section~\ref{sec:discussion}. Combined, the advantage holds over \textbf{four
city-pair directions} on three distinct cities, ruling out METR-LA-specific confounds.
Interestingly, GWNet[Strong] is \emph{worse} than default GWNet on PEMS04$\to$METR-LA
($12.58$ vs $11.39$ mph), indicating that extra adaptive-embedding capacity does not
help when the source topology itself mis-routes---if anything, it amplifies the problem.

\subsection{Cross-city zero-padding sensitivity}
\label{sec:appendix_no_pad_sensitivity}

The cross-city eval in
\texttt{experiments/stgnn\_shift\_eval\_v5\_pems04.py} (lines 87--91) uses
index-repetition padding when target $N$ < source $N$: for PEMS04 (307) $\to$
METR-LA (207) the \texttt{keep} list is the 207 unique target nodes followed
by 100 copies of node index 0, and MAE is averaged over all 307 prediction
slots. The duplicate slots feed a single repeated input, which can either
help (smoothing) or hurt (artificial concentration) the reported MAE. We
re-evaluate every PEMS04 checkpoint and recompute MAE only over the 207
\emph{unique} target nodes (no-pad), and compare to the existing 307-slot
padded MAE.

\begin{table}[h]
\centering
\small
\caption{Cross-city MAE (mph) on METR-LA target with the existing
index-repetition padding vs.\ a no-pad variant that scores only the 207
unique target nodes. PEMS04 $\to$ PEMS-BAY ($N{=}325$) is unaffected since
target $>$ source (subsample only). 3 PEMS04 seeds $\times$ 5 target subsets.
Script: \texttt{experiments/cross\_city\_no\_pad\_sensitivity.py};
data: \texttt{results/polish\_v15/cross\_city\_no\_pad\_sensitivity.json}.}
\label{tab:no_pad_sensitivity}
\vskip 0.05in
\begin{tabular}{lccc}
\toprule
\textbf{Backbone} & \textbf{Padded MAE (mph)} & \textbf{No-pad MAE (mph)} & \textbf{$\Delta$ (pad $-$ nopad)} \\
\midrule
DCRNN  & $7.014 \pm 0.085$  & $7.166 \pm 0.088$  & $-0.152$ \\
STGCN  & $7.176 \pm 0.169$  & $7.295 \pm 0.158$  & $-0.118$ \\
GWNet  & $11.390 \pm 0.623$ & $12.051 \pm 0.422$ & $-0.661$ \\
MTGNN  & $10.292 \pm 0.090$ & $11.020 \pm 0.083$ & $-0.728$ \\
GRU    & $11.026 \pm 0.008$ & $11.756 \pm 0.012$ & $-0.730$ \\
\bottomrule
\end{tabular}
\end{table}

The padding bias is small for the fixed-graph backbones DCRNN ($-0.15$ mph)
and STGCN ($-0.12$ mph) but larger for GWNet, MTGNN, and GRU
($-0.66$, $-0.73$, $-0.73$ mph): index-repetition padding flatters the
adaptive and graph-free models on PEMS04 $\to$ METR-LA by roughly $6$--$7\%$
of their already-large MAE. Crucially, the \emph{ranking} is unchanged: DCRNN
remains best in both columns, and the gap to the next-best model
(STGCN, $\approx 7$ vs.\ MTGNN, $\geq 10$) is preserved. The PEMS04 $\to$
PEMS-BAY direction is unaffected since target $>$ source (subsample only,
no padding).

% ============================================================
\section{Temporal Shift Analysis}
\label{sec:appendix_temporal_group}
% ============================================================

\subsection{Per-horizon decomposition (P19)}
\label{sec:appendix_horizon}

Table~\ref{tab:horizon} reports normalized MAE at horizons $h{=}1$, $h{=}6$, and
$h{=}12$ under clean, 30\% sensor-dropout, and cross-city (source-stat) evaluation
(3 seeds, mean).

\begin{table}[h]
\centering
\caption{Per-horizon normalized MAE on METR-LA at $h{=}1, 6, 12$.
Clean and 30\% sensor dropout computed on the METR-LA test set;
city-src computed on the 5-subset PEMS-BAY evaluation.
$\Delta$h: relative growth from $h{=}1$ to $h{=}12$.}
\label{tab:horizon}
\vskip 0.1in
\scriptsize
\setlength{\tabcolsep}{4pt}
\begin{tabular}{lccc|ccc|ccc}
\toprule
 & \multicolumn{3}{c}{\textbf{Clean}} & \multicolumn{3}{c}{\textbf{Masked 30\%}} & \multicolumn{3}{c}{\textbf{City-src}} \\
\cmidrule(lr){2-4}\cmidrule(lr){5-7}\cmidrule(lr){8-10}
\textbf{Backbone} & $h_1$ & $h_6$ & $h_{12}$ & $h_1$ & $h_6$ & $h_{12}$ & $h_1$ & $h_6$ & $h_{12}$ \\
\midrule
DCRNN         & $0.251$ & $0.356$ & $0.448$ & $0.356$ & $0.449$ & $0.527$ & $0.133$ & $0.160$ & $0.194$ \\
MTGNNLight    & $0.291$ & $0.370$ & $0.450$ & $0.371$ & $0.444$ & $0.515$ & $0.237$ & $0.240$ & $0.253$ \\
STGCN         & $0.296$ & $0.385$ & $0.466$ & $0.438$ & $0.513$ & $0.578$ & $0.173$ & $0.199$ & $0.230$ \\
GRU           & $0.331$ & $0.403$ & $0.474$ & $0.375$ & $0.435$ & $0.499$ & $0.263$ & $0.269$ & $0.282$ \\
GWNet         & $0.363$ & $0.423$ & $0.485$ & $0.464$ & $0.516$ & $0.568$ & $0.261$ & $0.264$ & $0.272$ \\
GWNet[Strong] & $0.333$ & $0.401$ & $0.468$ & $0.390$ & $0.457$ & $0.522$ & $0.254$ & $0.257$ & $0.260$ \\
\bottomrule
\end{tabular}
\end{table}

Under clean inference and sensor dropout, all backbones show the expected near-monotone
MAE growth with horizon. Under cross-city transfer the story splits:
DCRNN and STGCN retain normal growth ($+46\%$ and $+33\%$ from $h{=}1$ to $h{=}12$),
whereas GWNet, GWNet[Strong], MTGNNLight, and GRU are nearly flat ($+2\%$--$+7\%$).
For the adaptive-graph models, the transfer error is therefore paid \emph{up-front} as a
static structural mis-routing that saturates within the first prediction step---consistent
with the mask-aware adjacency intervention (Section~\ref{sec:mechanism}) recovering
75\% of the dropout gap without any temporal modification.

\subsection{Temporal covariate ablation}
\label{sec:appendix_temporal_cov}

To test whether the weekday/weekend gap of MegaCRN and STAEformer reflects an
over-reliance on their frozen time-of-day (TOD) and day-of-week (DOW)
embeddings, we run two inference-time interventions across five seeds per
model: (i) \emph{randomize}: replace TOD/DOW channels with uniform noise;
(ii) \emph{swap}: swap covariates between weekday and weekend samples.

\begin{table}[h]
\centering
\caption{Weekday~$\to$~weekend $\DeltaTemp$ (normalized) under covariate
interventions. Script: \texttt{experiments/temporal\_covariate\_ablation.py}.}
\label{tab:temporal_cov}
\vskip 0.05in
\small
\begin{tabular}{lccc|c}
\toprule
\textbf{Model} & \textbf{Clean} & \textbf{Randomize} & \textbf{Swap} & \textbf{Reduction (swap)} \\
\midrule
MegaCRN     & $0.7336$ & $0.7254$ & $0.7336$ & $0.0\%$ \\
STAEformer  & $0.7495$ & $0.7212$ & $0.6556$ & $12.5\%$ \\
\bottomrule
\end{tabular}
\end{table}

STAEformer shows partial covariate dependence; MegaCRN's gap is unchanged, so
its temporal fragility is not attributable to the frozen covariate channel.
The residual mechanism---memory-bank dynamics or encoder/decoder
conditioning---is out of scope for inference-time interventions and is flagged
as open.

% ============================================================
\section{Routing Concentration: Mechanism and Replication}
\label{sec:appendix_mechanism_group}
% ============================================================

\subsection{Adaptive adjacency density sweep}
\label{sec:appendix_density_sweep}

Table~\ref{tab:density_sweep} reports the anchor-point density sweep
($k \in \{10, 20, \text{full}\}$) alongside GWNet[fixed] and MTGNNLight references.
Table~\ref{tab:density_sweep_full} extends this to the full $k$-grid
$k \in \{5, 10, 15, 20, 30, \text{full}{=}207\}$.

\begin{table}[h]
\centering
\caption{Density sweep anchor points: relative fragility increases monotonically with
adaptive connectivity density.
Rel.\ fragility $= \DeltaSensor / \text{MAE}_\text{in-dist}$.
Mean $\pm$ std over 3 seeds; $\DeltaSensor$ averaged over 5 random 30\% masks.}
\label{tab:density_sweep}
\vskip 0.05in
\small
\begin{tabular}{lccc}
\toprule
\textbf{Variant} & \textbf{In-dist MAE} & \textbf{$\DeltaSensor$ (30\%)} & \textbf{Rel.\ fragility} \\
\midrule
GWNet[fixed]        & $8.27 \pm 0.08$ & $+0.050 \pm 0.010$ & $0.60\%$ \\
\midrule
GWNet[$k{=}10$]     & $9.03 \pm 0.42$ & $+0.073 \pm 0.030$ & $0.80\%$ \\
GWNet[$k{=}20$]     & $8.85 \pm 0.34$ & $+0.074 \pm 0.032$ & $0.84\%$ \\
GWNet[$k{=}\text{full}$] & $7.85 \pm 0.11$ & $+0.092 \pm 0.030$ & $1.17\%$ \\
\midrule
MTGNNLight (ref)    & $6.88 \pm 0.11$ & $+0.051 \pm 0.012$ & $0.74\%$ \\
\bottomrule
\end{tabular}
\end{table}
All evaluations use frozen GWNet[adaptive] checkpoints (3 seeds);
top-$k$ filtering is applied at inference only, with row renormalization.
Sanity check: $k{=}\text{full}$ reproduces registry values within $10^{-5}$ MAE (verified).

\begin{table}[h]
\centering
\caption{Full $k$-grid density sweep for GWNet[adaptive] on METR-LA.
Mean $\pm$ std over 3 seeds.
$\DeltaSensor$: averaged over 5 random 30\% zero-masks (seeds 42--46).
Rel.\ fragility $= \DeltaSensor / \text{MAE}_\text{in-dist}$.
Note: $k{<}\text{full}$ disrupts embeddings trained for dense routing, causing
in-distribution degradation at low $k$; the clean-performance cost at $k{=}5$
is $+17\%$ relative to $k{=}\text{full}$.}
\label{tab:density_sweep_full}
\vskip 0.1in
\small
\begin{tabular}{lcccc}
\toprule
\textbf{$k$} & \textbf{In-dist MAE} & \textbf{$\DeltaSensor$ (30\%)} & \textbf{Rel.\ fragility} \\
\midrule
$5$          & $9.21 \pm 0.50$ & $+0.072 \pm 0.031$ & $0.78\%$ \\
$10$         & $9.03 \pm 0.42$ & $+0.073 \pm 0.030$ & $0.80\%$ \\
$15$         & $8.94 \pm 0.38$ & $+0.073 \pm 0.031$ & $0.81\%$ \\
$20$         & $8.85 \pm 0.34$ & $+0.074 \pm 0.032$ & $0.84\%$ \\
$30$         & $8.66 \pm 0.30$ & $+0.078 \pm 0.033$ & $0.89\%$ \\
$\text{full}=207$ & $7.85 \pm 0.11$ & $+0.092 \pm 0.030$ & $1.17\%$ \\
\bottomrule
\end{tabular}
\end{table}

The relative fragility trend is monotonically increasing across all six $k$ values
(5/5 consecutive pairs), supporting the association between denser adaptive connectivity
and higher dropout fragility within this checkpoint.
The absolute $\DeltaSensor$ values at $k{=}5$--$20$ are similar ($0.072$--$0.074$),
with the gap widening at $k{=}30$--$\text{full}$.
As noted in Section~\ref{sec:routing_spectrum}, inference-time sparsification is not
a practical intervention because it degrades in-distribution accuracy; the finding
is presented as mechanism evidence, not as a deployment recommendation.

\subsection{Routing-sparsity spectrum across adaptive models}
\label{sec:appendix_routing_spectrum}

Table~\ref{tab:routing_spectrum} reports the per-model adjacency concentration
and sensor-dropout delta for the four adaptive-graph models in our bank
(METR-LA test set; three seeds). $\text{max\_mass}$ is the mean over rows and
samples of the largest entry in each adjacency row; $\text{eff\_}k$ is
$\exp(\text{row-entropy})$ and measures effective support size.

STAEformer is excluded from the Pearson because its spatial mixing is
head-stratified multi-head attention rather than a single row-stochastic
adjacency: a max\_mass extracted from any one head is not commensurable
with the per-row probability used in the other four backbones, so the
proxy-isomorphism assumption underlying the correlation does not hold.
The four-backbone Pearson is therefore reported on the row-stochastic
subset where the proxy is identically defined.

\begin{table}[h]
\centering
\caption{Routing-sparsity spectrum vs.\ sensor-dropout fragility.
Pearson correlation of $\text{max\_mass}$ and $\DeltaSensor$ across four adaptive
models: $\rho = {+}0.979$. Script:
\texttt{experiments/routing\_sparsity\_spectrum.py}. $\DeltaSensor$ values are
the routing-spectrum estimates used for this correlation.}
\label{tab:routing_spectrum}
\vskip 0.05in
\small
\begin{tabular}{lccccc}
\toprule
\textbf{Model} & \textbf{Family} & \textbf{max\_mass} & \textbf{eff\_$k$} & \textbf{density} & \textbf{$\DeltaSensor$} \\
\midrule
AGCRN        & dense-symmetric  & $0.008$ & $\approx 200$ & $1.00$ & ${+}0.011$ \\
MegaCRN      & memory-meta      & $0.008$ & $\approx 200$ & $1.00$ & ${-}0.005$ \\
MTGNN-Light  & sparse-top-$k$   & $0.160$ & ---           & sparse & ${+}0.051$ \\
GWNet        & dense-bilinear   & $0.230$ & $\approx 114$ & $1.00$ & ${+}0.092$ \\
\bottomrule
\end{tabular}
\end{table}

Under the conservative ranking that treats MegaCRN's negative delta as an
adjacency-smoothing artifact (Section~\ref{sec:polish_ablations}), max\_mass
and $\DeltaSensor$ are monotonic across all four models.

\subsubsection{Seed-level variation and rank-correlation diagnostics}
\label{sec:appendix_routing_spectrum_seedvar}

Per-seed max\_mass values for the four adaptive backbones are AGCRN
$\{0.0077,\,0.0076,\,0.0100\}$ (std $0.0014$), MegaCRN
$\{0.0064,\,0.0093,\,0.0097\}$ (std $0.0018$), MTGNN-Light
$\{0.1586,\,0.1607,\,0.1596\}$ (std $0.0011$), and GWNet
$\{0.1638,\,0.2773,\,0.2475\}$ (std $0.0588$). Three of the four backbones
have per-seed std below the smallest non-tied inter-family gap of $0.07$
(MTGNN $0.16$ vs.\ GWNet $0.23$); GWNet's larger seed spread is the dominant
source of within-family variability.

Per-seed Pearson against $\DeltaSensor$ is $\{+0.92,\,+0.99,\,+0.98\}$ across
the three seeds (mean $+0.96$, std $0.04$), suggesting that the max\_mass
relation is not an artifact of seed averaging despite GWNet's larger
within-seed spread. Spearman over the four
seed-mean values is $\rho_S{=}{+}0.80$; the single rank exchange swaps AGCRN
($0.0085$) and MegaCRN ($0.0085$), whose max\_mass values differ by less than
$10^{-3}$, well below either backbone's per-seed std. Treating the two
memory-meta backbones as tied on the proxy yields rank correlation
$\rho_S{=}{+}1.0$ across the three distinguishable groups
$\{\text{tied AGCRN/MegaCRN}, \text{MTGNN-Light}, \text{GWNet}\}$.

\subsection{PEMS04-native routing-concentration replication}
\label{sec:appendix_pems04_routing}

To address whether the routing-concentration relation generalizes beyond a single source
city, we re-train the four adaptive-graph backbones on PEMS04 (307 sensors,
chronological 70/10/20 split) using the same training configuration as the
canonical METR-LA bank (Section~\ref{subsec:bank}). $\text{max\_mass}$ is then
recomputed from each PEMS04-trained checkpoint via the same per-architecture
adjacency proxy (Section~\ref{sec:routing_spectrum}), and $\DeltaSensor$ is
measured on PEMS04 test data (3 model seeds $\times$ 5 mask seeds, 30\%
zero-mask, MAE on kept sensors only).

\begin{table}[h]
\centering
\caption{Routing-concentration relation replicates on PEMS04. Pearson
$\rho(\text{max\_mass}, \DeltaSensor) = {+}0.910$ across the four adaptive
backbones (METR-LA reference: ${+}0.979$). Native checkpoints; same protocol
as Table~\ref{tab:routing_spectrum}. Script:
\texttt{experiments/pems04\_routing\_gate\_native.py};
data: \texttt{results/polish\_v15/pems04\_native\_routing\_gate.json}.}
\label{tab:pems04_routing}
\vskip 0.05in
\small
\begin{tabular}{lcccc}
\toprule
& \multicolumn{2}{c}{\textbf{METR-LA (reference)}} & \multicolumn{2}{c}{\textbf{PEMS04 (native)}} \\
\cmidrule(lr){2-3}\cmidrule(lr){4-5}
\textbf{Model} & \textbf{max\_mass} & \textbf{$\DeltaSensor$} & \textbf{max\_mass} & \textbf{$\DeltaSensor$} \\
\midrule
AGCRN     & $0.008$ & $+0.011$ & $0.006 \pm 0.000$ & $+0.051 \pm 0.006$ \\
MegaCRN   & $0.008$ & $-0.005$ & $0.005 \pm 0.000$ & $+0.057 \pm 0.018$ \\
MTGNN     & $0.160$ & $+0.051$ & $0.145 \pm 0.001$ & $+0.119 \pm 0.004$ \\
GWNet     & $0.230$ & $+0.092$ & $0.248 \pm 0.020$ & $+0.114 \pm 0.025$ \\
\bottomrule
\end{tabular}
\end{table}

Two findings strengthen the shared-mechanism hypothesis. (i) max\_mass values are tightly preserved
across source cities (within a few percent for each backbone), confirming
that adjacency concentration is an architecture-level property rather than a
dataset artifact. (ii) MegaCRN's METR-LA-specific negative $\DeltaSensor$
($-0.005$) becomes positive on PEMS04 ($+0.057$), consistent with the
adjacency-mediated weekday-prior smoothing diagnosis in
Section~\ref{sec:polish_ablations}---the smoothing prior exists in METR-LA
but not in the PEMS04 weekday distribution, while the routing-concentration
mechanism is shared.

\subsection{Theory sketch: first-order exposure bound}
\label{sec:appendix_theory}

\paragraph{Setup.}
Let a single message-passing layer produce node representations
$h_i = \sum_j A_{ij} W x_j$ where $A \in \mathbb{R}^{N\times N}$ is
row-stochastic and $W$ is a linear map. Perturbing the inputs at nodes in
$\mathcal{I} \subseteq [N]$ by zero-masking, the change in $h_i$ is
$\Delta h_i = -\sum_{j \in \mathcal{I}} A_{ij} W x_j$.

\paragraph{Exposure bound.}
Taking norms and bounding input magnitudes by $\|x\|_\infty$:
\begin{equation*}
  \|\Delta h_i\|_2 \leq \|W\|_2 \cdot \|x\|_\infty \cdot \sum_{j \in \mathcal{I}} A_{ij}.
\end{equation*}
For any row $i$, let $m_i=\max_j A_{ij}$. Then
$\sum_{j \in \mathcal{I}} A_{ij} \leq \lvert \mathcal{I} \rvert m_i$, so

\begin{equation*}
  \|\Delta h_i\|_2 \leq \|W\|_2 \cdot \|x\|_\infty \cdot
  \lvert \mathcal{I} \rvert m_i.
\end{equation*}
Averaging over target nodes yields
\begin{equation*}
  \frac{1}{N}\sum_i \|\Delta h_i\|_2
  \leq \|W\|_2 \cdot \|x\|_\infty \cdot
  \lvert \mathcal{I} \rvert \cdot \text{max\_mass}(A).
\end{equation*}

\paragraph{Interpretation.}
The bound says that concentrated rows have larger worst-case exposure when
the dropped set includes high-mass neighbors. It is deliberately local: it
does not prove full-horizon MAE behavior and does not cover nonlinear
recurrence, temporal convolutions, or decoder amplification. It also should
not be read as an exact expected-error result under iid dropout. If each
sensor is dropped independently with probability $p$, then for a row-stochastic
$A$,
\begin{equation*}
  \mathbb{E}_{\mathcal{I}}\left[\sum_{j\in\mathcal{I}} A_{ij}\right]
  = p\sum_j A_{ij}=p,
\end{equation*}
which is independent of $m_i$. Thus $\text{max\_mass}(A)$ is best viewed as a
tail-risk or worst-case exposure proxy. This is the role it plays in
Table~\ref{tab:routing_spectrum}: the empirical $\rho{=}{+}0.979$ on METR-LA
and $\rho{=}{+}0.910$ on PEMS04 show that the proxy tracks sensor-dropout
fragility across adaptive backbones. The mask-aware intervention
(Section~\ref{sec:mechanism}) targets the same exposure term by zeroing and
renormalizing routes attached to known dropped sensors.

% ============================================================
\section{Model Ablations}
\label{sec:appendix_ablations_group}
% ============================================================

\subsection{GWNet capacity variants (P17 and P21)}
\label{sec:appendix_stronger_gwnet}

We train GWNet[Strong] with hidden dim 128, adaptive embed dim 20, and early-stopping
patience 30 (otherwise identical to the standard config). Table~\ref{tab:stronger_gwnet}
compares it to the default GWNet and DCRNN across all shift axes.
We additionally report GWNet[Best] (P21), the best configuration identified by
in-distribution MAE across a 12-config Stage-1 grid
(hidden $\in\{128,256\}$, embed $\in\{20,32\}$, lr $\in\{1\text{e-}3,5\text{e-}4,2\text{e-}4\}$,
dropout $\in\{0.1,0.3\}$); the winning setting is hidden 256, embed 32, lr 5e-4, patience 30.

\begin{table}[h]
\centering
\caption{GWNet capacity scaling: standard, P17 strong, and P21 best-sweep configs vs.\ DCRNN
(3 seeds, mean $\pm$ std).
Source-only adjacency MAE follows the P16 fairness control protocol.
Q4/Q1 is the degree-stratified fragility ratio (Appendix~\ref{sec:appendix_degree_stratification}).
GWNet[Best] city MAE (mph) and Q4/Q1 were not separately evaluated in the sweep;
source-only normalized values are reported.}
\label{tab:stronger_gwnet}
\vskip 0.1in
\small
\begin{tabular}{lccccc}
\toprule
\textbf{Variant} & \textbf{In-dist} & \textbf{$\DeltaSensor$} & \textbf{City MAE} & \textbf{Source-only} & \textbf{Q4/Q1} \\
 & \textbf{(mph)} & \textbf{(30\%)} & \textbf{(mph)} & \textbf{adj (norm)} & \textbf{ratio} \\
\midrule
DCRNN         & $6.58 \pm 0.03$ & $+0.020 \pm 0.002$ & $2.74 \pm 0.01$ & $0.180 \pm 0.001$ & $0.95\times$ \\
GWNet         & $7.85 \pm 0.11$ & $+0.088 \pm 0.017$ & $4.42 \pm 0.15$ & $0.268 \pm 0.006$ & $1.83\times$ \\
GWNet[Strong] & $7.42 \pm 0.04$ & $+0.050 \pm 0.018$ & $4.46 \pm 0.06$ & $0.262 \pm 0.004$ & $1.72\times$ \\
GWNet[Best]   & $7.22 \pm 0.04$ & $+0.042 \pm 0.002$ & ---             & $0.262 \pm 0.009$ & ---          \\
\bottomrule
\end{tabular}
\end{table}

Doubling hidden dimension and adaptive embedding closes 34\% of the in-distribution MAE
gap to DCRNN ($7.85 \to 7.42$ mph) and halves the sensor-dropout $\Delta$ ($+0.088 \to +0.050$),
matching the sparse MTGNNLight baseline. However, the cross-city MAE is effectively
unchanged ($4.42 \to 4.46$ mph), the source-only adjacency gap to DCRNN persists
($0.262$ vs $0.180$), and the degree-stratified hub-fragility signature remains
monotonic ($1.72\times$ Q4/Q1, still far from DCRNN's $0.95\times$).

The full 12-config sweep GWNet[Best] (hidden 256, embed 32, lr 5e-4, patience 30)
further reduces the in-distribution gap ($7.22 \pm 0.04$ mph, closing 63\% of the gap
to DCRNN) and the sensor-dropout fragility ($+0.042 \pm 0.002$), yet
$\DeltaSensor$ remains $2.1\times$ DCRNN's $+0.020$---confirming P21.
The source-only city-transfer normalized MAE ($0.262 \pm 0.009$) is identical to
GWNet[Strong]'s ($0.262$) and still 46\% above DCRNN's $0.180$.
Capacity attenuates but does not eliminate the architectural fragility, supporting the
mechanism interpretation in Section~\ref{sec:mechanism}.

\subsection{MegaCRN memory-bank ablation}
\label{sec:appendix_memory_ablation}

To test whether MegaCRN's meta-memory bank absorbs missing sensors as a form of
implicit regularization, we zero-mask and re-route its attention/memory paths
at inference in four modes. \emph{full}: the unmodified model; \emph{memzero}:
query the memory bank with a zero vector so the prototype fall-back is fixed
regardless of input; \emph{graphonly}: zero out the attention output keeping
the adjacency pathway; \emph{attnonly}: use identity adjacency while keeping
memory attention. We report $\DeltaSensor$ means over three seeds and five
dropout ratios $r \in \{0.1, 0.2, 0.3, 0.4, 0.5\}$.

\begin{table}[h]
\centering
\caption{MegaCRN $\DeltaSensor$ (normalized) by dropout ratio and ablation
mode. Script: \texttt{experiments/megacrn\_memory\_ablation.py}.}
\label{tab:memory_ablation}
\vskip 0.05in
\small
\begin{tabular}{lccccc}
\toprule
\textbf{Mode}    & $r{=}0.1$ & $r{=}0.2$ & $r{=}0.3$ & $r{=}0.4$ & $r{=}0.5$ \\
\midrule
full        & ${-}0.001$ & ${-}0.043$ & ${-}0.040$ & ${-}0.016$ & ${+}0.045$ \\
memzero     & ${-}0.006$ & ${-}0.067$ & ${-}0.075$ & ${-}0.038$ & ${-}0.007$ \\
graphonly   & ${+}0.001$ & ${-}0.027$ & ${-}0.024$ & ${+}0.012$ & ${+}0.068$ \\
attnonly    & $0.000$    & $0.000$    & $0.000$    & $0.000$    & $0.000$    \\
\bottomrule
\end{tabular}
\end{table}

\textbf{memzero} deepens rather than flattens the U-shape, refuting the
memory-bank regularization hypothesis. \textbf{attnonly} (identity adjacency)
eliminates $\DeltaSensor$ completely, establishing that cross-node flow
through the adjacency is the sole driver. Per-seed $\sigma$ at $r{=}0.2$ is
$0.031$, approximately equal to the mean---the apparent robustness is
seed-fragile. The mechanism is adjacency-mediated smoothing: zero-diffused
predictions drift toward a weekday-mean prior that accidentally tracks the
test set.

% ============================================================
\section{Architecture-Selection Decision Tree}
\label{sec:appendix_decision_tree}
% ============================================================

Table~\ref{tab:decision_tree} reports the conservative $2{\times}2{\times}2$
recommendation table over the deployment profile. The tree is generated by
\texttt{experiments/decision\_tree\_artifact.py} from the same per-model
normalized fragility vectors used in the main paper. The conservative ranking
demotes flagged \emph{(model, axis)} pairs where low apparent fragility is an
artifact (\emph{(megacrn, sensor)}; see Section~\ref{sec:polish_ablations}).

\begin{table}[h]
\centering
\caption{Architecture recommendation per deployment profile. Margin is the
difference in weighted-risk score to the runner-up (lower = closer call).
Confidence is margin normalized by runner-up risk.}
\label{tab:decision_tree}
\vskip 0.05in
\small
\begin{tabular}{cccccc}
\toprule
\textbf{sensor} & \textbf{city} & \textbf{weekend} & \textbf{Recommended} & \textbf{Runner-up} & \textbf{Confidence} \\
\midrule
no  & no  & no  & DCRNN      & AGCRN      & $0.30$ \\
no  & no  & yes & AGCRN      & MTGNN-Light& $1.00$ \\
no  & yes & no  & MegaCRN    & DCRNN      & $1.00$ \\
no  & yes & yes & DCRNN      & AGCRN      & $0.51$ \\
yes & no  & no  & STAEformer & AGCRN      & $0.18$ \\
yes & no  & yes & AGCRN      & DCRNN      & $0.19$ \\
yes & yes & no  & DCRNN      & MegaCRN    & $0.34$ \\
yes & yes & yes & DCRNN      & AGCRN      & $0.30$ \\
\bottomrule
\end{tabular}
\end{table}

Four of the eight cells reach DCRNN under multi-shift risk, reflecting DCRNN's
uniformly low fragility across axes. MegaCRN is recommended only under the
single-shift city-transfer profile; STAEformer wins the sensor-only cell;
AGCRN wins the two weekend-exposed cells without city transfer. The tree is
released as a pip-installable anonymous artifact (\texttt{stgnn-stress}) so
that practitioners can reproduce the recommendation and the underlying
rankings on new frozen checkpoints.
